%% ============================================================
%%  Response to Reviewers — HGST Preprint v1 → v2
%%  "Frustration Universality in HGST: U(1) to the SM gauge group"
%%  Date: March 5, 2026
%% ============================================================
\documentclass[12pt,a4paper]{article}
\usepackage[T1]{fontenc}
\usepackage[utf8]{inputenc}
\usepackage{lmodern}
\usepackage{geometry}
\usepackage{xcolor}
\usepackage{hyperref}
\usepackage{enumitem}
\usepackage{mdframed}
\usepackage{amsmath,amssymb}
\usepackage{setspace}
\geometry{top=2.5cm,bottom=2.5cm,left=3cm,right=3cm}
\onehalfspacing

\definecolor{authorblue}{RGB}{0,70,140}
\definecolor{reviewergray}{RGB}{80,80,80}
\definecolor{changedgreen}{RGB}{0,120,50}

\newmdenv[
  linecolor=gray!50,
  backgroundcolor=gray!6,
  roundcorner=3pt,
  innertopmargin=8pt,
  innerbottommargin=8pt,
  skipabove=6pt,skipbelow=6pt
]{reviewerbox}

\newmdenv[
  linecolor=authorblue!60,
  backgroundcolor=authorblue!5,
  roundcorner=3pt,
  innertopmargin=8pt,
  innerbottommargin=8pt,
  skipabove=6pt,skipbelow=6pt
]{authorbox}

\newcommand{\reviewer}[1]{\textbf{\textcolor{reviewergray}{Reviewer #1}}}
\newcommand{\RC}[2]{\noindent\begin{reviewerbox}\textbf{\textcolor{reviewergray}{RC\,#1.#2:}}}
\newcommand{\RCend}{\end{reviewerbox}}
\newcommand{\AR}[2]{\noindent\begin{authorbox}\textbf{\textcolor{authorblue}{AR\,#1.#2:}}}
\newcommand{\ARend}{\end{authorbox}}
\newcommand{\changed}[1]{\textcolor{changedgreen}{\textit{[Changed: #1]}}}

%% ============================================================
\begin{document}
%% ============================================================

\begin{center}
  {\Large\bfseries Response to Referees}\\[8pt]
  {\large Frustration Universality in Hierarchical Graded Structure Theory:\\
  A Lattice Study of the E7 MIXED Class from U(1) to the Standard Model Gauge Group}\\[6pt]
  Manuscript revised \today \\[4pt]
  {\small Authors: HGST Development Team}
\end{center}

\bigskip
\noindent\rule{\linewidth}{0.8pt}

\noindent
We thank the referees for their careful reading and constructive comments.
Below we address every point raised.
Changes to the manuscript are highlighted in blue in the revised PDF
(produced with \texttt{latexdiff main\_v1.tex main.tex}).
All numerical results are unchanged; changes are editorial/clarificatory
unless noted otherwise.

\bigskip
\noindent\rule{\linewidth}{0.4pt}

%% ============================================================
\section*{Referee 1}
%% ============================================================

%% --- Template: copy/paste the RC+AR block pair for each comment ---

\RC{1}{1}
  \textit{[Paste referee comment here.]}
\RCend

\AR{1}{1}
  We thank the referee for this observation. [Response here.]
  \changed{Section X, paragraph Y.}
\ARend

\RC{1}{2}
  \textit{[Next comment.]}
\RCend

\AR{1}{2}
  [Response here.]
\ARend

%% Add more RC+AR pairs as needed for Referee 1

\bigskip
\noindent\rule{\linewidth}{0.4pt}

%% ============================================================
\section*{Referee 2}
%% ============================================================

\RC{2}{1}
  \textit{[Paste referee comment here.]}
\RCend

\AR{2}{1}
  [Response here.]
\ARend

%% ============================================================
\section*{List of All Changes (v1 $\to$ v2)}
%% ============================================================

\noindent The following table summarises every substantive change made to the manuscript.

\medskip
\begin{enumerate}[leftmargin=1.5em]
  %% Fill in as revisions are made
  \item \textbf{[§X, line Y]:} [Description of change.]
\end{enumerate}

\end{document}
